\clearpage
\thispagestyle{plain}
\phantomsection
\addcontentsline{toc}{chapter}{KATA PENGANTAR}

\begin{center}
{\fontfamily{ptm}\selectfont\bfseries KATA PENGANTAR}
\end{center}
\vspace{0.5cm}

% Isi Konten
{\fontfamily{ptm}\selectfont
Puji dan syukur kepada Tuhan Yang Maha Esa, atas berkat dan rahmatNya sehingga penulis dapat menyelesaikan laporan Skripsi ini dengan baik dan tepat waktu. Penulis juga tidak lupa berterimakasih kepada pihak yang telah membantu penulis dalam memberikan bimbingan, dukungan, dan doa. Oleh karena itu, pada kesempatan ini penulis ingin mengucapkan terima kasih kepada:
\begin{enumerate}
    \item Ibu Dr. Nelly, S.Kom., M.M., selaku Rektor Binus University.
    \item Bapak Dr. Ir. Derwin Suhartono, S.Kom., MTI, selaku Dekan of School of Computer Science and Head of Computer Science Binus Bandung.
    \item Bapak Budi Juarto, S.T., M.Kom., selaku dosen pembimbing yang telah banyak meluangkan waktu untuk memberikan bimbingan, petunjuk, dukungan, dan saran kepada penulis dalam menyelesaikan laporan Skripsi ini.
    \item Segenap Dosen Fakultas School of Computer Science yang telah mendidik dan memberikan ilmu selama kuliah dan seluruh staf yang selalu sabar melayani segala administrasi selama proses penelitian ini.
    \item Orang tua dan saudara-saudari yang senantiasa mendampingi dengan doa, semangat, dan dukungan materi, sehingga penulis mampu melalui proses pendidikan dan menyelesaikan skripsi ini dengan baik.
    \item Semua pihak yang telah membantu dan mendukung penulis selama pembuatan laporan Skripsi ini yang tidak dapat penulis sebutkan satu per satu.
\end{enumerate}

% Paragraf Penutup
Karena keterbatasan pengetahuan maupun pengalaman penulis, maka penulis yakin masih banyak kekurangan dalam laporan Skripsi ini, Oleh karena itu penulis sangat mengharapkan saran dan kritik yang membangun dari pembaca demi kesempurnaan dokumen skripsi ini.

\vspace{1.5cm}
 
\begin{flushright}
Bandung, 16 Desember 2025\\[2cm] 
\textbf{Penulis}
\end{flushright}
}

\newpage